% section3_2_frequentist.tex — Module A: Yue
\subsection{Frequentist Baseline Comparison}\label{sec:frequentist}

This section presents the ordinary least squares (OLS) estimates and compares them with the Bayesian posterior summaries from Section~\ref{sec:posterior}.

\subsubsection{OLS Parameter Estimates}

Table~\ref{tab:freq_vs_bayes} compares the maximum likelihood estimates (MLE) with posterior means and their respective 95\% intervals.

\begin{table}[htbp]
\centering
\caption{Frequentist MLE vs.\ Bayesian posterior estimates for $\texttt{cnt} \sim \texttt{temp} + \texttt{hum}$.}
\label{tab:freq_vs_bayes}
\begin{tabular}{lcccccc}
\hline
Parameter & MLE & 95\% CI & Post.\ Mean & 95\% CrI \\
\hline
$\beta_0$ (intercept) & TODO & TODO & TODO & TODO \\
$\beta_1$ (temp)       & TODO & TODO & TODO & TODO \\
$\beta_2$ (hum)        & TODO & TODO & TODO & TODO \\
$\sigma$               & TODO & --- & TODO & TODO \\
\hline
\end{tabular}
\end{table}

\subsubsection{Model Diagnostics}

% TODO: fill from frequentist_diagnostics.csv
\begin{itemize}
    \item $R^2 = \text{[r\_squared]}$, Adjusted $R^2 = \text{[adj\_r\_squared]}$
    \item $F$-statistic $= \text{[f\_statistic]}$ ($p < \text{[f\_pvalue]}$)
    \item RMSE $= \text{[rmse]}$
    \item Breusch--Pagan test for heteroscedasticity: LM statistic $= \text{[bp\_lm\_stat]}$, $p = \text{[bp\_lm\_pvalue]}$
\end{itemize}

% TODO: interpret BP test result
If the Breusch--Pagan $p$-value is small (e.g., $< 0.05$), this suggests heteroscedasticity in the residuals, which would affect the reliability of standard errors under OLS but does not invalidate the Bayesian model if the posterior predictive checks account for this.

\subsubsection{Interval Interpretation}

The frequentist 95\% confidence interval and the Bayesian 95\% credible interval have fundamentally different interpretations:

\begin{itemize}
    \item \textbf{Confidence interval (CI)}: Under repeated sampling, 95\% of such intervals would contain the true parameter value. The specific interval either does or does not contain the truth --- we cannot assign a probability to this event.
    \item \textbf{Credible interval (CrI)}: Given the observed data and prior, there is a 0.95 posterior probability that the parameter lies within this interval. This is a direct probability statement about the parameter.
\end{itemize}

In practice, when the prior is weakly informative and the sample size is moderate ($n = 360$), the CI and CrI are expected to be numerically similar, reflecting the dominance of the likelihood. Differences arise primarily from:
\begin{enumerate}
    \item The regularising effect of the prior (shrinkage toward prior mean);
    \item The Bayesian treatment of $\sigma$ as a random variable rather than a plug-in estimate.
\end{enumerate}

\subsubsection{Summary}

The frequentist baseline confirms that temperature is a strong positive predictor of daily rentals, while humidity has a weaker negative association. These findings are consistent with the Bayesian posterior and provide a useful reference point for evaluating the added value of the Bayesian framework (uncertainty quantification, sequential updating, prior sensitivity).
